% GENERATED FILE. Edit canonical lessons, then run npm run generate:books.
% canonical-source-hash: fnv1a64:d9bbf7398e4b170a
% canonical-lessons: LA-C489-laetitia, LA-C489-laetor, LA-C489-hilaritas, LA-C489-exsulto, LA-C489-iucunditas

\chapter{Joy and gladness}
\label{ch:la-joy-and-gladness}
\hlchaptermodality{489}

\begin{chapteropening}
\textbf{By the end of this chapter:} \emph{I can talk about joy and gladness.}

Complete the last lesson of chapter 489: I can talk about joy and gladness.
\end{chapteropening}

\section[laetitia, laetitiae]{laetitia, laetitiae — joy, gladness}

\label{lesson:LA-C489-laetitia}



\begin{quote}
Before the new one: say the Latin for a ditch, then the Latin for a rampart.
\end{quote}

\subsection*{You'll want to know: laetitia, laetitiae}
\textbf{laetitia, laetitiae} — \textquotedblleft{}joy, gladness\textquotedblright{}.

\begin{quote}
Post victōriam magna laetitia in urbe erat.
\end{quote}

\begin{quote}
\emph{After the victory there was great joy in the city.}
\end{quote}

Feminine, first declension. \emph{Laetitia} is lively joy that shows; \emph{gaudium} is joy felt inside.

\begin{cousinweb}
\textbf{laetitia} is built on the adjective \emph{laetus} \textquotedblleft{}glad, happy\textquotedblright{} with the noun ending \emph{-itia} (Lewis \& Short; EtymDB derives it from \emph{laetus}).
\end{cousinweb}

\begin{grammarlens}[title={What you've built}]
One more word for telling, explaining and describing at B1.
\end{grammarlens}

\subsection*{Your turn}
\begin{itemize}
\raggedright
  \item \emph{Say it:} \emph{laetitia, laetitiae}
  \item \emph{Say it:} \emph{laetitia, laetitiae}, once more
  \item \emph{Say it:} a sentence of your own with \emph{laetitia, laetitiae}
  \item \emph{Recall:} say the Latin for a ditch, then the Latin for a rampart, then say \emph{laetitia, laetitiae} again
\end{itemize}

\begin{tcolorbox}[breakable,colback=teal!4,colframe=teal!35!black,title={Before you move on}]
What does laetitia, laetitiae mean? (\textquotedblleft{}Joy, gladness\textquotedblright{}.) Say it once more.
\end{tcolorbox}
\section[laetor, laetārī]{laetor, laetārī — to rejoice, be glad}

\label{lesson:LA-C489-laetor}



\begin{quote}
Before the new one: say the Latin for a rampart, then the Latin for joy.
\end{quote}

\subsection*{You'll want to know: laetor, laetārī}
\textbf{laetor, laetārī} — \textquotedblleft{}to rejoice, be glad\textquotedblright{}.

\begin{quote}
Māter laetātur, quod fīlius sānus redit.
\end{quote}

\begin{quote}
\emph{The mother rejoices because her son comes home healthy.}
\end{quote}

Deponent: \emph{laetor, laetārī, laetātus sum}. With the ablative or with \emph{dē}: \emph{victōriā laetor} — I rejoice at the victory.

\begin{cousinweb}
\textbf{laetor} is a deponent verb built on \emph{laetus} \textquotedblleft{}glad, happy\textquotedblright{} (Lewis \& Short).
\end{cousinweb}

\begin{grammarlens}[title={What you've built}]
One more word for telling, explaining and describing at B1.
\end{grammarlens}

\subsection*{Your turn}
\begin{itemize}
\raggedright
  \item \emph{Say it:} \emph{laetor, laetārī}
  \item \emph{Say it:} \emph{laetor, laetārī}, once more
  \item \emph{Say it:} a sentence of your own with \emph{laetor, laetārī}
  \item \emph{Recall:} say the Latin for a rampart, then the Latin for joy, then say \emph{laetor, laetārī} again
\end{itemize}

\begin{tcolorbox}[breakable,colback=teal!4,colframe=teal!35!black,title={Before you move on}]
What does laetor, laetārī mean? (\textquotedblleft{}To rejoice, be glad\textquotedblright{}.) Say it once more.
\end{tcolorbox}
\section[hilaritās, hilaritātis]{hilaritās, hilaritātis — cheerfulness, merriment}

\label{lesson:LA-C489-hilaritas}



\begin{quote}
Before the new one: say the Latin for joy, then the Latin for to rejoice.
\end{quote}

\subsection*{You'll want to know: hilaritās, hilaritātis}
\textbf{hilaritās, hilaritātis} — \textquotedblleft{}cheerfulness, merriment\textquotedblright{}.

\begin{quote}
Hilaritās avī omnēs hospitēs laetōs facit.
\end{quote}

\begin{quote}
\emph{Grandfather's cheerfulness makes all the guests happy.}
\end{quote}

Feminine, third declension. English \emph{hilarity} comes from it.

\begin{cousinweb}
\textbf{hilaritās} is built on \emph{hilaris} \textquotedblleft{}cheerful\textquotedblright{}, an early borrowing from Greek \emph{hilaros} \textquotedblleft{}cheerful\textquotedblright{} (EtymDB; Lewis \& Short). English \emph{hilarity} comes from it.
\end{cousinweb}

\begin{grammarlens}[title={What you've built}]
One more word for telling, explaining and describing at B1.
\end{grammarlens}

\subsection*{Your turn}
\begin{itemize}
\raggedright
  \item \emph{Say it:} \emph{hilaritās, hilaritātis}
  \item \emph{Say it:} \emph{hilaritās, hilaritātis}, once more
  \item \emph{Say it:} a sentence of your own with \emph{hilaritās, hilaritātis}
  \item \emph{Recall:} say the Latin for joy, then the Latin for to rejoice, then say \emph{hilaritās, hilaritātis} again
\end{itemize}

\begin{tcolorbox}[breakable,colback=teal!4,colframe=teal!35!black,title={Before you move on}]
What does hilaritās, hilaritātis mean? (\textquotedblleft{}Cheerfulness, merriment\textquotedblright{}.) Say it once more.
\end{tcolorbox}
\section[exsultō, exsultāre]{exsultō, exsultāre — to leap for joy, rejoice greatly}

\label{lesson:LA-C489-exsulto}



\begin{quote}
Before the new one: say the Latin for to rejoice, then the Latin for cheerfulness.
\end{quote}

\subsection*{You'll want to know: exsultō, exsultāre}
\textbf{exsultō, exsultāre} — \textquotedblleft{}to leap for joy, rejoice greatly\textquotedblright{}.

\begin{quote}
Puerī in hortō exsultant, quod hodiē fēriae sunt.
\end{quote}

\begin{quote}
\emph{The boys leap for joy in the garden, because today is a holiday.}
\end{quote}

\emph{exsultō, exsultāre, exsultāvī, exsultātum}. First to leap up, then to be overjoyed. English \emph{exult}.

\begin{cousinweb}
\textbf{exsultō} is built on \emph{ex-} \textquotedblleft{}out, up\textquotedblright{} and \emph{saltō} \textquotedblleft{}to jump, dance\textquotedblright{}, the frequentative of \emph{saliō} \textquotedblleft{}to leap\textquotedblright{} (Lewis \& Short). English \emph{exult} comes from it.
\end{cousinweb}

\begin{grammarlens}[title={What you've built}]
One more word for telling, explaining and describing at B1.
\end{grammarlens}

\subsection*{Your turn}
\begin{itemize}
\raggedright
  \item \emph{Say it:} \emph{exsultō, exsultāre}
  \item \emph{Say it:} \emph{exsultō, exsultāre}, once more
  \item \emph{Say it:} a sentence of your own with \emph{exsultō, exsultāre}
  \item \emph{Recall:} say the Latin for to rejoice, then the Latin for cheerfulness, then say \emph{exsultō, exsultāre} again
\end{itemize}

\begin{tcolorbox}[breakable,colback=teal!4,colframe=teal!35!black,title={Before you move on}]
What does exsultō, exsultāre mean? (\textquotedblleft{}To leap for joy, rejoice greatly\textquotedblright{}.) Say it once more.
\end{tcolorbox}
\section[iūcunditās, iūcunditātis]{iūcunditās, iūcunditātis — pleasantness, delight}

\label{lesson:LA-C489-iucunditas}



\begin{quote}
Before the new one: say the Latin for cheerfulness, then the Latin for to leap for joy.
\end{quote}

\subsection*{You'll want to know: iūcunditās, iūcunditātis}
\textbf{iūcunditās, iūcunditātis} — \textquotedblleft{}pleasantness, delight\textquotedblright{}.

\begin{quote}
Iūcunditātem hortī nostrī omnēs amīcī laudant.
\end{quote}

\begin{quote}
\emph{All our friends praise the pleasantness of our garden.}
\end{quote}

Feminine, third declension, built on \emph{iūcundus} — pleasant.

\begin{cousinweb}
\textbf{iūcunditās} is built on \emph{iūcundus} \textquotedblleft{}pleasant, agreeable\textquotedblright{} with the noun ending \emph{-itās} (Lewis \& Short, under \emph{jucunditas}); Lewis \& Short links \emph{iūcundus} with \emph{iocus} \textquotedblleft{}a joke\textquotedblright{}, but its origin is not settled.
\end{cousinweb}

\begin{grammarlens}[title={What you've built}]
One more word for telling, explaining and describing at B1.
\end{grammarlens}

\subsection*{Your turn}
\begin{itemize}
\raggedright
  \item \emph{Say it:} \emph{iūcunditās, iūcunditātis}
  \item \emph{Say it:} \emph{iūcunditās, iūcunditātis}, once more
  \item \emph{Say it:} a sentence of your own with \emph{iūcunditās, iūcunditātis}
  \item \emph{Recall:} say the Latin for cheerfulness, then the Latin for to leap for joy, then say \emph{iūcunditās, iūcunditātis} again
\end{itemize}

\begin{tcolorbox}[breakable,colback=teal!4,colframe=teal!35!black,title={Before you move on}]
What does iūcunditās, iūcunditātis mean? (\textquotedblleft{}Pleasantness, delight\textquotedblright{}.) Say it once more.
\end{tcolorbox}
